%% ============================================================
%%  Aegis: Distributed Host-Based Intrusion Detection System
%%  Mini Project II Report — WCE Sangli, TY CSE 2025-26
%%  Template adapted from Mini Project I (Light Automation)
%% ============================================================

\documentclass[12pt,a4paper]{report}

%% ----- Packages -----
\usepackage[T1]{fontenc}
\usepackage[utf8]{inputenc}
\usepackage{geometry}
\usepackage{graphicx}
\usepackage{fancyhdr}
\usepackage{titlesec}
\usepackage{booktabs}
\usepackage{longtable}
\usepackage{tabularx}
\usepackage{array}
\usepackage{multirow}
\usepackage{hyperref}
\usepackage{xcolor}
\usepackage{listings}
\usepackage{enumitem}
\usepackage{setspace}
\usepackage{tocloft}
\usepackage{caption}
\usepackage{float}
\usepackage{amsmath}
\usepackage{parskip}
\usepackage{mdframed}
\usepackage{tikz}
\usepackage{eso-pic}
\usepackage{lmodern}

%% ----- Page Border (All Pages) -----
\AddToShipoutPictureBG{%
  \AtPageLowerLeft{%
    \begin{tikzpicture}[remember picture,overlay]
      \draw[line width=1pt]
        ($(current page.north west)+(1cm,-1cm)$) rectangle
        ($(current page.south east)+(-1cm,1cm)$);
    \end{tikzpicture}%
  }%
}

%% ----- Page Geometry -----
\geometry{
  a4paper,
  top=1in, bottom=1in,
  left=1.25in, right=1in
}

%% ----- Hyperref Setup -----
\hypersetup{
  colorlinks=true,
  linkcolor=black,
  urlcolor=blue,
  citecolor=black
}

%% ----- Header/Footer -----
\pagestyle{fancy}
\fancyhf{}
\fancyhead[L]{\small Aegis: Distributed HIDS}
\fancyhead[R]{\small WCE Sangli, TY CSE 2025-26}
\fancyfoot[C]{\thepage}
\renewcommand{\headrulewidth}{0.4pt}

%% ----- Code Listing Style -----
\definecolor{codebg}{RGB}{245,245,245}
\definecolor{codecomment}{RGB}{100,100,100}
\definecolor{codestring}{RGB}{0,128,0}
\definecolor{codekeyword}{RGB}{0,0,180}

\lstdefinestyle{javastyle}{
  backgroundcolor=\color{codebg},
  commentstyle=\color{codecomment}\itshape,
  keywordstyle=\color{codekeyword}\bfseries,
  stringstyle=\color{codestring},
  basicstyle=\ttfamily\small,
  breakatwhitespace=false,
  breaklines=true,
  keepspaces=true,
  numbers=left,
  numberstyle=\tiny\color{gray},
  showspaces=false,
  showstringspaces=false,
  tabsize=2,
  frame=single,
  rulecolor=\color{gray!50}
}

\lstdefinestyle{propstyle}{
  backgroundcolor=\color{codebg},
  basicstyle=\ttfamily\small,
  breaklines=true,
  frame=single,
  rulecolor=\color{gray!50}
}

%% ----- Section Title Formatting -----
\titleformat{\chapter}[block]
  {\normalfont\Large\bfseries}{\thechapter.}{1em}{}
\titlespacing*{\chapter}{0pt}{20pt}{10pt}

%% ----- Spacing -----
\setstretch{1.3}

%% ============================================================
\begin{document}

%% ============================================================
%%  TITLE PAGE
%% ============================================================
\begin{titlepage}
  \thispagestyle{empty}
  \centering

  % College logo — place wcelogo.png in the same directory
  \includegraphics[width=3cm]{wcelogo.png}\\[10pt]

  {\LARGE\bfseries Walchand College of Engineering, Sangli}\\[4pt]
  {\large (Government aided Autonomous Institute)}\\[8pt]
  \rule{\linewidth}{0.5pt}\\[8pt]
  {\large\bfseries Department of Computer Science and Engineering}\\[6pt]
  {\normalsize TY CSE Mini Project II}\\[4pt]
  {\normalsize Report on}\\[16pt]

  {\LARGE\bfseries Aegis: Distributed Host-Based Intrusion Detection System}\\[20pt]

  {\normalsize Submitted by}\\[10pt]
  \begin{tabular}{ll}
    \toprule
    \textbf{Name} & \textbf{PRN} \\
    \midrule
    Avadhut Bharat Mali      & 23510003 \\
    Yash Shitalkumar Savalkar & 23510027 \\
    Aditya Rajkumar Pansare   & 23510080 \\
    \bottomrule
  \end{tabular}\\[20pt]

  {\normalsize Under the guidance of}\\[6pt]
  {\large\bfseries Dr.\ A.\ R.\ Surve}\\[4pt]
  {\normalsize Department of Computer Science and Engineering,}\\
  {\normalsize Walchand College of Engineering, Sangli}\\[20pt]

  {\large\bfseries 2025 -- 2026}
\end{titlepage}

%% ============================================================
%%  CERTIFICATE
%% ============================================================
\chapter*{Certificate}
\addcontentsline{toc}{chapter}{Certificate}
\thispagestyle{plain}

\begin{center}
  {\large\bfseries Walchand College of Engineering, Sangli}\\
  {\normalsize (Government aided Autonomous Institute)}\\[6pt]
  {\large\bfseries Department of Computer Science and Engineering}\\[16pt]
  {\LARGE\bfseries\color{blue} CERTIFICATE}
\end{center}

\vspace{10pt}

This is to certify that the Project Report entitled, \textbf{AEGIS: DISTRIBUTED HOST-BASED INTRUSION DETECTION SYSTEM} submitted by Mr.\ Avadhut Bharat Mali, Mr.\ Yash Shitalkumar Savalkar, and Mr.\ Aditya Rajkumar Pansare, to Walchand College of Engineering Sangli, India, is a record of bonafide project work of course Mini Project II carried out by them under our supervision and guidance and is worthy of consideration for the award of the degree of Bachelor of Technology in Computer Science \& Engineering of the Institute.

\vspace{200pt}

\begin{center}
\begin{tabular}{p{3.5cm} p{4.5cm} p{3.5cm} p{3.5cm}}

\centering \textbf{Dr.\ A.\ R.\ Surve} &
\centering \textbf{Mr.\ H.\ R.\ Chandoba} &
\centering \textbf{External Examiner} &
\centering \textbf{Dr.\ A.\ R.\ Surve} \tabularnewline[6pt]

\centering Guide &
\centering Mini Project Coordinator &
&
\centering Head of Department \tabularnewline[10pt]

\centering Dept.\ of CSE &
\centering Dept.\ of CSE &
&
\centering Dept.\ of CSE \tabularnewline

\centering WCE Sangli &
\centering WCE Sangli &
&
\centering WCE Sangli

\end{tabular}
\end{center}
%% ============================================================
%%  ACKNOWLEDGEMENT
%% ============================================================
\chapter*{Acknowledgement}
\addcontentsline{toc}{chapter}{Acknowledgement}

We would like to express our sincere gratitude to all those who guided and supported us throughout the successful completion of our Mini Project titled \textbf{``AEGIS: DISTRIBUTED HOST-BASED INTRUSION DETECTION SYSTEM''}.

We are deeply thankful to \textbf{Dr.\ A.\ R.\ Surve}, our project guide, for his constant encouragement, valuable suggestions, and expert guidance at every stage of this project. His mentorship has been instrumental in helping us design a multi-layered security system and in navigating the complexities of distributed architecture.

We extend our heartfelt thanks to \textbf{Dr.\ A.\ R.\ Surve}, Head of the Department of Computer Science and Engineering, for providing us with the necessary facilities and an environment conducive to research and innovation.

We are also grateful to the Mini Project Coordinator and all the faculty members of the Department of Computer Science and Engineering, Walchand College of Engineering, Sangli, for their continuous motivation, support, and constructive feedback during the development of our project.

%% ============================================================
%%  DECLARATION
%% ============================================================
\chapter*{Declaration}
\addcontentsline{toc}{chapter}{Declaration}

We hereby declare that the work presented in this project report titled \textbf{AEGIS: DISTRIBUTED HOST-BASED INTRUSION DETECTION SYSTEM} submitted by us in the partial fulfillment of the requirement for the award of the degree of Bachelor of Technology (B.Tech) in the Department of Computer Science \& Engineering, Walchand College of Engineering, Sangli, is an authentic record of our project work carried out under the guidance of Dr.\ A.\ R.\ Surve.

\vspace{20pt}

Date: \underline{\hspace{3cm}} \hfill Place: Sangli

\vspace{20pt}

\begin{tabular}{ll}
  Avadhut Bharat Mali       & PRN: 23510003 \\[10pt]
  Yash Shitalkumar Savalkar & PRN: 23510027 \\[10pt]
  Aditya Rajkumar Pansare   & PRN: 23510080 \\
\end{tabular}

%% ============================================================
%%  TABLE OF CONTENTS
%% ============================================================
\vspace{300pt}
\tableofcontents
\listoffigures
\listoftables

%% ============================================================
%%  ACRONYMS
%% ============================================================
\chapter*{Acronyms}
\addcontentsline{toc}{chapter}{Acronyms}

\begin{tabular}{ll}
  \textbf{HIDS}     & Host-based Intrusion Detection System \\
  \textbf{NIDS}     & Network-based Intrusion Detection System \\
  \textbf{BYOD}     & Bring Your Own Device \\
  \textbf{ERP}      & Enterprise Resource Planning \\
  \textbf{BSSID}    & Basic Service Set Identifier (router MAC address) \\
  \textbf{SSID}     & Service Set Identifier (Wi-Fi network name) \\
  \textbf{FIM}      & File Integrity Monitoring \\
  \textbf{JWT}      & JSON Web Token \\
  \textbf{REST}     & Representational State Transfer \\
  \textbf{AMQP}     & Advanced Message Queuing Protocol \\
  \textbf{API}      & Application Programming Interface \\
  \textbf{SPA}      & Single Page Application \\
  \textbf{ORM}      & Object Relational Mapping \\
  \textbf{CI/CD}    & Continuous Integration / Continuous Deployment \\
  \textbf{STOMP}    & Simple Text Oriented Messaging Protocol \\
  \textbf{DPDP}     & Digital Personal Data Protection (Act 2023) \\
  \textbf{SaaS}     & Software as a Service \\
  \textbf{SOC}      & Security Operations Centre \\
  \textbf{SIEM}     & Security Information and Event Management \\
  \textbf{JNA}      & Java Native Access \\
  \textbf{AES}      & Advanced Encryption Standard \\
  \textbf{HTTPS}    & Hypertext Transfer Protocol Secure \\
  \textbf{HSTS}     & HTTP Strict Transport Security \\
  \textbf{DLQ}      & Dead Letter Queue \\
  \textbf{LAN}      & Local Area Network \\
  \textbf{PRN}      & Permanent Registration Number \\
  \textbf{TTL}      & Time To Live \\
  \textbf{UUID}     & Universally Unique Identifier \\
\end{tabular}

%% ============================================================
\setcounter{chapter}{0}

%% ============================================================
%%  CHAPTER 1 — ABSTRACT
%% ============================================================
\chapter{Abstract}

Aegis is a custom-designed Distributed Host-Based Intrusion Detection System (HIDS) developed to secure the Bring Your Own Device (BYOD) ecosystem at Walchand College of Engineering, Sangli. Traditional perimeter defenses—such as the existing Sophos firewall—are often insufficient against internal endpoint vulnerabilities and localized campus threats because they provide only network-level visibility while remaining blind to host-level behaviour.

Aegis addresses this security gap by deploying four coordinated layers: a lightweight \textbf{Android Agent} (Kotlin) on student smartphones, a \textbf{PC Agent} (Java 17) daemon on laboratory computers, a \textbf{Spring Boot 3 backend server} with RabbitMQ message brokering, and a \textbf{React 18 admin dashboard} for real-time campus-wide monitoring.

The system serves a dual purpose: it enhances user convenience through automated Sophos portal login and secure ERP access, while simultaneously enforcing a Zero-Trust security model by evaluating real-time device health via a calculated Security Score. If a device is compromised, Aegis dynamically restricts its access to critical college infrastructure. The distributed agents actively detect Evil Twin (Rogue Wi-Fi) access points, unauthorized APK installations, file integrity violations, suspicious processes, and malicious network connections. All telemetry flows through RabbitMQ to a centralized analysis engine that correlates events, identifies Patient Zero for malware outbreaks, and pushes live alerts to the admin dashboard via WebSocket.

By combining automated campus services with crowdsourced endpoint monitoring, Aegis provides a robust, real-time security layer that protects both individual students and the broader college network without compromising personal privacy, in full compliance with India's Digital Personal Data Protection Act 2023.

%% ============================================================
%%  CHAPTER 2 — INTRODUCTION AND RELATED WORK
%% ============================================================
\chapter{Introduction and Related Work}

\section{Background}

The Walchand College network relies primarily on centralized perimeter defenses (Sophos firewalls) that are effective for external threats but leave individual student devices exposed. The WCE ERP is protected only by CAPTCHA, and the result portal uses OTP which can be intercepted. Sophos sees traffic; it cannot see what generates that traffic on a device. This creates a critical blind spot between the network perimeter and device-level behaviour—the exact gap that Aegis is designed to fill.

Modern campus environments are characterized by the BYOD model: hundreds of student smartphones and laboratory PCs connect to the same network simultaneously. A single compromised endpoint can become a Patient Zero for lateral malware spread across the LAN.

\section{Evolution of Intrusion Detection Systems}

A Host-Based Intrusion Detection System (HIDS) monitors and analyzes the internal activities of a computing system, providing deeper visibility than Network-based IDSs (NIDS). While NIDS watches traffic moving across the network, HIDS focuses on host-level events like file changes, registry modifications, and suspicious processes.

A major limitation identified in traditional HIDS implementations is that they can create a performance bottleneck on the host machine if not tuned correctly, leading to resource starvation on the monitored endpoint.

\section{Analysis of Existing Solutions}

\textbf{OSSEC} has long been a pioneer in the open-source HIDS space, utilizing a traditional client-server model to analyze system activities. However, it has historically struggled with scalability when deployed across large environments and lacks modern API support for seamless integrations.

\textbf{Wazuh}, which forked from OSSEC in 2015, addressed these limitations by improving scalability, offering continuous updates, and integrating robust event indexing. Despite these advancements, full-scale enterprise tools like Wazuh or the ELK stack can be highly complex and resource-heavy for small-to-medium environments—identifying a clear need for lightweight, highly concurrent alternatives.

\section{File Integrity Monitoring Challenges}

File Integrity Monitoring (FIM) is a cornerstone of endpoint forensics that establishes a baseline of file attributes using cryptographic hashes (commonly SHA-256). The system continuously checks files against this baseline to detect unauthorized changes. A significant challenge is the trade-off between real-time detection speed and system performance—real-time monitoring of high-activity directories can generate large alert volumes leading to ``alert fatigue'' for security teams.

\section{Problem Identification and Project Justification}

Existing enterprise HIDS solutions are often monolithic, demanding significant host resources to process and format logs locally before transmission. To solve this, Aegis introduces a strictly decoupled architecture: by offloading event queuing to RabbitMQ and delegating all heavy analytical processing to a centralized Spring Boot backend, the host agent remains exceptionally lightweight. This ensures that continuous forensic data gathering does not bottleneck the monitored system.

%% ============================================================
%%  CHAPTER 3 — PROBLEM STATEMENT
%% ============================================================
\chapter{Problem Statement}

The Walchand College network relies primarily on centralized perimeter defenses (Sophos firewalls). While effective for external threats, this approach leaves individual student endpoints vulnerable and creates several critical security gaps:

\begin{itemize}[leftmargin=*, noitemsep]
  \item \textbf{Manual Authentication Risks:} Daily manual logins to the Sophos portal (\texttt{http://192.168.1.8:8090}) and Walchand ERP increase user fatigue, making students susceptible to credential theft via Evil Twin (Rogue Wi-Fi) and phishing attacks.

  \item \textbf{Lack of Endpoint Visibility:} Existing network firewalls cannot detect host-level threats such as toxic application permissions, rooted OS vulnerabilities, or malicious background processes on student devices.

  \item \textbf{Unrestricted Infrastructure Access:} Any device connected to the network can access critical infrastructure like the Walchand ERP without any Zero-Trust health verification, risking lateral malware spread.

  \item \textbf{No Centralized Threat Monitoring:} College IT administrators lack a distributed, real-time dashboard to track localized campus threats (e.g., a fake access point in the canteen) or quickly isolate compromised devices.

  \item \textbf{No PC-level Forensics:} Laboratory computers have no mechanism to detect unauthorized file system changes (FIM), suspicious process execution, or malicious outbound network connections.
\end{itemize}

%% ============================================================
%%  CHAPTER 4 — OBJECTIVES
%% ============================================================
\chapter{Objectives}

The primary objective is to develop Aegis, a comprehensive and privacy-preserving Distributed Host-Based Intrusion Detection System for the Walchand College campus. Specifically, the project aims:

\begin{enumerate}[leftmargin=*, noitemsep]
  \item \textbf{Automate secure authentication:} Seamlessly authenticate users to the Sophos network portal and Walchand ERP, reducing user friction and mitigating credential theft via phishing.

  \item \textbf{Assess real-time endpoint health:} Develop lightweight host agents (Android \& PC) that continuously calculate a Security Score by monitoring OS integrity (root/jailbreak detection) and identifying toxic application permissions.

  \item \textbf{Implement a Zero-Trust security model:} Dynamically restrict a device's access to critical college infrastructure (ERP) if its Security Score falls below an acceptable threshold.

  \item \textbf{Detect localized network threats:} Identify Evil Twin (Rogue Wi-Fi) access points and malicious captive portals at the endpoint level before a connection is exploited.

  \item \textbf{Establish File Integrity Monitoring:} Monitor critical system directories on laboratory PCs using SHA-256 hash baselines to detect unauthorized file modifications, creations, or deletions.

  \item \textbf{Monitor processes and network connections:} Flag suspicious process spawning and outbound connections to known malicious IPs on PC agents.

  \item \textbf{Build a centralized threat dashboard:} Provide IT administrators with a real-time panel for campus-wide security event monitoring, device health boards, and Patient Zero identification.

  \item \textbf{Ensure DPDP Act 2023 compliance:} Collect only security-posture metadata (never personal content) and obtain explicit user consent before enrollment.
\end{enumerate}

%% ============================================================
%%  CHAPTER 5 — SYSTEM ARCHITECTURE
%% ============================================================
\chapter{System Architecture}

\section{Overview}

Aegis employs a four-layer distributed client-server architecture. The two lightweight agents (Android and PC) collect host-level telemetry and report to a Spring Boot backend via HTTPS REST APIs. The backend queues incoming events through RabbitMQ, processes them asynchronously via a rule engine, stores results in PostgreSQL, caches live state in Redis, and pushes real-time alerts to the admin dashboard via WebSocket (STOMP over SockJS).

\begin{figure}[H]
  \centering
  % Replace with your actual architecture diagram image
  % \includegraphics[width=\linewidth]{architecture_flow.png}
  \fbox{\parbox{0.9\linewidth}{\centering\vspace{40pt}
    \textit{[Insert Architecture Flow Diagram here]}\\
    Android Agent $\rightarrow$ PC Agent $\rightarrow$ Spring Boot API $\rightarrow$ RabbitMQ $\rightarrow$ Analysis Engine $\rightarrow$ Dashboard
    \vspace{40pt}}}
  \caption{Aegis System Architecture Flow}
  \label{fig:arch}
\end{figure}

\section{The Four Layers}

\begin{description}[leftmargin=*, noitemsep]
  \item[Layer 1 — Android Agent (Kotlin):] Runs on student smartphones. Handles Evil Twin detection, Security Score calculation, Sophos auto-login, ERP gated access, and threat reporting.

  \item[Layer 2 — PC Agent (Java 17):] Runs as a daemon on laboratory computers. Performs File Integrity Monitoring (FIM), process monitoring, and network connection scanning.

  \item[Layer 3 — Backend Server (Spring Boot 3):] Ingests all telemetry via REST, queues events through RabbitMQ, scores devices via a rule engine, enforces ERP access policies, and pushes live alerts via WebSocket.

  \item[Layer 4 — Admin Dashboard (React 18):] Provides a real-time threat map, device health board, alert feed, Patient Zero tracker, analytics, and policy manager to IT administrators.
\end{description}

\section{Key Design Decisions}

\begin{itemize}[leftmargin=*, noitemsep]
  \item Agents \textbf{never connect directly to PostgreSQL}. All writes go through REST API $\rightarrow$ RabbitMQ $\rightarrow$ consumer. This decouples ingestion from processing and handles the 9am burst (200+ devices connecting simultaneously) without database contention.

  \item \textbf{RabbitMQ} with three priority queues ensures high-severity threat alerts are processed before low-priority heartbeats.

  \item \textbf{Redis} caches device state (current score, online/offline) for sub-millisecond dashboard lookups without hitting PostgreSQL on every WebSocket refresh.

  \item All external communication goes through \textbf{nginx} on port 443 only—all other services (PostgreSQL, RabbitMQ, Redis) live on an internal Docker bridge network.
\end{itemize}

%% ============================================================
%%  CHAPTER 6 — TECHNOLOGY STACK
%% ============================================================
\chapter{Technology Stack}

\section{Android Agent}

\begin{longtable}{@{}p{2.5cm}p{3cm}p{2cm}p{6cm}@{}}
  \toprule
  \textbf{Layer} & \textbf{Technology} & \textbf{Version} & \textbf{Purpose / Rationale} \\
  \midrule
  \endhead
  Language    & Kotlin          & 1.9+      & Coroutines for async, null safety, minimal boilerplate \\
  UI          & Jetpack Compose & 1.5+      & Modern declarative UI for enrollment, score, alert screens \\
  Background  & WorkManager     & 2.9+      & Survives Android 12+ battery optimization that kills services \\
  Local DB    & Room            & 2.6+      & Typed SQLite; offline event buffer before server sync \\
  HTTP        & Retrofit + OkHttp & 2.9/4.12 & REST calls with certificate pinning via OkHttp interceptor \\
  Crypto      & Android Keystore & API 23+  & Hardware-backed AES-256-GCM; keys never leave secure enclave \\
  \bottomrule
  \caption{Android Agent Technology Stack}
\end{longtable}

\section{PC Agent}

\begin{longtable}{@{}p{2.5cm}p{3cm}p{2cm}p{6cm}@{}}
  \toprule
  \textbf{Layer} & \textbf{Technology} & \textbf{Version} & \textbf{Purpose / Rationale} \\
  \midrule
  \endhead
  Language    & Java 17         & LTS       & Cross-platform daemon; same language as backend; shared DTOs \\
  Build       & Maven           & 3.9+      & Fat JAR for single-file distribution \\
  File Monitor & WatchService   & JDK built-in & Pure Java FIM; cross-platform; no native dependency \\
  OS Calls    & JNA             & 5.14+     & Process enumeration; simpler than JNI; maps C structs directly \\
  Local DB    & H2 Embedded     & 2.2+      & FIM hash baseline; zero install; SQL diff queries \\
  HTTP        & OkHttp (Java)   & 4.12      & Retry + exponential backoff built-in \\
  \bottomrule
  \caption{PC Agent Technology Stack}
\end{longtable}

\section{Backend Server}

\begin{longtable}{@{}p{2.5cm}p{3.5cm}p{2cm}p{5.5cm}@{}}
  \toprule
  \textbf{Layer} & \textbf{Technology} & \textbf{Version} & \textbf{Purpose / Rationale} \\
  \midrule
  \endhead
  Framework   & Spring Boot 3   & 3.2+      & REST API + WebSocket; team's strongest backend skill \\
  Auth        & Spring Security + JWT & 6.x / 0.12 & Stateless JWT for agents; refresh rotation for admin \\
  Message Broker & RabbitMQ    & 3.12+     & Decouples ingestion; handles burst; dead-letter queue \\
  ORM         & Spring Data JPA & 3.2+      & Repository pattern; typed queries \\
  Primary DB  & PostgreSQL      & 15+       & Structured schema; JOINs for analytics; JSONB for payloads \\
  Cache       & Redis           & 7.x       & Sub-ms score lookup; session tracking; 24h TTL \\
  Migrations  & Flyway          & 9.x       & Reproducible DB setup; no manual SQL on fresh deploy \\
  WebSocket   & Spring WebSocket + STOMP & 6.x & Live dashboard push; integrates with SockJS \\
  \bottomrule
  \caption{Backend Server Technology Stack}
\end{longtable}

\section{Admin Dashboard}

\begin{longtable}{@{}p{2.5cm}p{3.5cm}p{2cm}p{5.5cm}@{}}
  \toprule
  \textbf{Layer} & \textbf{Technology} & \textbf{Version} & \textbf{Purpose / Rationale} \\
  \midrule
  \endhead
  Framework   & React 18        & 18.x      & Component reuse; concurrent rendering for live feed \\
  Language    & TypeScript      & 5.x       & Catches WebSocket payload type errors at compile time \\
  Styling     & Tailwind + shadcn/ui & 3.4   & Prebuilt DataTable, Badge, Dialog; zero bundle overhead \\
  WebSocket   & SockJS + STOMP.js & latest  & SockJS fallback to polling if WS blocked on campus network \\
  Maps        & Leaflet.js      & 1.9+      & Free; no API key; works offline on campus intranet \\
  Charts      & Chart.js        & 4.x       & Score trends + alert volume bar charts \\
  Data Fetch  & TanStack Query  & 5.x       & Background refetch; loading states built-in \\
  State       & Zustand         & 4.x       & Auth, active alerts, WebSocket connection stores \\
  \bottomrule
  \caption{Admin Dashboard Technology Stack}
\end{longtable}

%% ============================================================
%%  CHAPTER 7 — METHODOLOGY
%% ============================================================
\chapter{Methodology}

Aegis employs a four-phase operational methodology across its distributed components.

\section{Phase 1: Endpoint Environmental Scanning (Host Agents)}

\textbf{Android — Evil Twin Detection:} A \texttt{BroadcastReceiver} fires on \texttt{NETWORK\_STATE\_CHANGED}. Upon Wi-Fi connection, the agent fetches the current BSSID via \texttt{WifiManager} and cross-checks it against a server-maintained list of legitimate WCE MAC addresses. A CRITICAL alert is fired if a rogue AP is detected; the user is notified and the event is reported to the backend.

\textbf{Android — Security Score Engine:} Calculates a 0–100 score using the formula: start at 100, deduct $-30$ for root, $-20$ for developer options enabled, $-15$ for screen lock disabled, up to $-25$ for apps with dangerous permissions (not whitelisted), $-5$ for unknown BSSID. Runs via WorkManager every 30 minutes.

\textbf{PC — File Integrity Monitoring:} On first run, generates SHA-256 baseline hashes of watched directories (\texttt{/etc}, \texttt{/bin}, \texttt{/usr/bin} on Linux; \texttt{C:\textbackslash Windows\textbackslash System32} on Windows) stored in H2. \texttt{WatchService} detects CREATE/MODIFY/DELETE events and diffs against the baseline. Smart debouncing batches events within a 5-second window.

\textbf{PC — Process Monitor:} \texttt{ProcessHandle.allProcesses()} snapshots every 60 seconds. Flags new processes matching suspicious patterns: unexpected \texttt{powershell.exe}, \texttt{cmd.exe} spawned by non-user, \texttt{netcat}, Python one-liners. JNA checks process parent PID for privilege escalation on Windows.

\textbf{PC — Network Scanner:} Parses \texttt{netstat} output every 30 seconds. Flags connections to IPs in a locally cached AbuseIPDB blocklist (refreshed daily) and port-scan behavior (>10 unique destination IPs in 60 seconds).

\section{Phase 2: Secure Automation and Zero-Trust Enforcement}

\textbf{Sophos Auto-Login:} On Wi-Fi connect event, OkHttp POST to \texttt{192.168.1.8:8090/userLogin} with Keystore-encrypted credentials. No manual portal interaction required.

\textbf{ERP Gating:} A WebView with \texttt{FLAG\_SECURE} (blocks screen capture) is gated by the device's current Security Score. If the score is below the server-configurable threshold (default $\geq 60$), access is denied with a reason screen. ERP credentials auto-fill from Android Keystore.

\section{Phase 3: Distributed Threat Reporting}

When an anomaly is detected, the agent packages it as a \texttt{ThreatEventDTO} JSON payload and POSTs it to \texttt{/api/v1/threats/report} with JWT authentication. On network failure, events are stored in the local Room DB (Android) or H2 (PC) and retried with exponential backoff (1s, 2s, 4s, max 60s). The backend responds with HTTP 202 (Accepted) immediately and queues the event in RabbitMQ for asynchronous processing.

\section{Phase 4: Centralized Management (Admin Dashboard)}

The backend aggregates all incoming telemetry, correlates events across devices for Patient Zero detection (time-window clustering), and pushes real-time alerts via WebSocket STOMP topics to the React dashboard. The dashboard presents a live threat map, device health board, alert feed, analytics charts, and a policy manager through which administrators can update ERP score thresholds, manage known BSSID lists, and control app whitelists—changes propagated to all agents on their next heartbeat.

%% ============================================================
%%  CHAPTER 8 — PROJECT DIAGRAMS
%% ============================================================
\chapter{Project Diagrams}

\section{Functional Block Diagram}

\begin{figure}[H]
  \centering
  % \includegraphics[width=0.85\linewidth]{functional_block.png}
  \fbox{\parbox{0.85\linewidth}{\centering\vspace{60pt}
    \textit{[Insert Functional Block Diagram — Student Connects to WCE\_WIFI $\rightarrow$ Agent Checks BSSID $\rightarrow$ Mismatch (Evil Twin) / Match (Legit Wi-Fi) $\rightarrow$ Security Score $\rightarrow$ ERP gating / Sophos auto-login]}
    \vspace{60pt}}}
  \caption{Aegis Functional Block Diagram}
  \label{fig:fbd}
\end{figure}

\section{UML Sequence Diagram}

\begin{figure}[H]
  \centering
  % \includegraphics[width=0.9\linewidth]{uml_sequence.png}
  \fbox{\parbox{0.9\linewidth}{\centering\vspace{60pt}
    \textit{[Insert UML Sequence Diagram — Aegis Mobile Agent $\leftrightarrow$ Sophos Portal $\leftrightarrow$ Walchand ERP $\leftrightarrow$ Central Admin Server]}
    \vspace{60pt}}}
  \caption{UML Sequence Diagram — Agent to Server Interaction}
  \label{fig:uml}
\end{figure}

\section{Database Entity Relationship Diagram}

The database schema consists of five primary tables: \textbf{devices}, \textbf{security\_events}, \textbf{heartbeats}, \textbf{access\_policies}, and \textbf{patient\_zero\_clusters}.

\begin{figure}[H]
  \centering
  \fbox{\parbox{0.9\linewidth}{\centering\vspace{60pt}
    \textit{[Insert ER Diagram — devices $\leftarrow$ security\_events, heartbeats; access\_policies; patient\_zero\_clusters]}
    \vspace{60pt}}}
  \caption{Database Entity Relationship Diagram}
  \label{fig:erd}
\end{figure}

\subsection*{Key Database Tables}

\begin{longtable}{@{}p{3.5cm}p{3cm}p{2cm}p{5cm}@{}}
  \toprule
  \textbf{Column} & \textbf{Type} & \textbf{Constraint} & \textbf{Description} \\
  \midrule
  \multicolumn{4}{l}{\textbf{Table: devices}} \\
  \midrule
  id              & UUID          & PK          & Device UUID, generated on enrollment \\
  prn             & VARCHAR(20)   & NOT NULL    & Student PRN (pseudonymous link) \\
  device\_type    & ENUM          & NOT NULL    & ANDROID \textbar{} PC \\
  current\_score  & INT           & DEFAULT 100 & Latest security score 0--100 \\
  risk\_level     & ENUM          & NOT NULL    & CLEAN \textbar{} SUSPICIOUS \textbar{} COMPROMISED \\
  erp\_access\_blocked & BOOLEAN & DEFAULT false & Admin or auto block flag \\
  consent\_accepted & BOOLEAN    & NOT NULL    & DPDP Act 2023 consent flag \\
  \midrule
  \multicolumn{4}{l}{\textbf{Table: security\_events}} \\
  \midrule
  id              & BIGSERIAL     & PK          & Auto-increment event ID \\
  device\_id      & UUID          & FK devices  & Source device \\
  event\_type     & ENUM          & NOT NULL    & EVIL\_TWIN, FIM\_CHANGE, PROC\_ANOMALY, etc. \\
  severity        & ENUM          & NOT NULL    & CRITICAL \textbar{} HIGH \textbar{} MEDIUM \textbar{} LOW \\
  payload         & JSONB         &             & Flexible event data (BSSID, path, IP, etc.) \\
  acknowledged    & BOOLEAN       & DEFAULT false & Admin acknowledged flag \\
  \bottomrule
  \caption{Selected Database Schema Fields}
\end{longtable}

%% ============================================================
%%  CHAPTER 9 — API DESIGN
%% ============================================================
\chapter{API Design}

All API endpoints are served at base URL \texttt{https://aegis.wce.ac.in/api/v1}. All endpoints require \texttt{Authorization: Bearer <JWT>} unless marked PUBLIC.

\section{Agent Endpoints}

\begin{longtable}{@{}p{1.5cm}p{5cm}p{2cm}p{5cm}@{}}
  \toprule
  \textbf{Method} & \textbf{Path} & \textbf{Auth} & \textbf{Description} \\
  \midrule
  \endhead
  POST & /agent/register    & PUBLIC     & Device enrollment; returns deviceId, JWT, policy \\
  POST & /agent/heartbeat   & Agent JWT  & Periodic score + telemetry push; returns policy updates \\
  POST & /threats/report    & Agent JWT  & Report threat event; queued in RabbitMQ; returns 202 \\
  GET  & /agent/policy      & Agent JWT  & Pull current policy (BSSID list, whitelist, score threshold) \\
  POST & /agent/token/refresh & PUBLIC   & Refresh expired JWT using refresh token \\
  \bottomrule
  \caption{Agent REST Endpoints}
\end{longtable}

\section{Admin Endpoints}

\begin{longtable}{@{}p{1.5cm}p{6cm}p{2cm}p{4cm}@{}}
  \toprule
  \textbf{Method} & \textbf{Path} & \textbf{Auth} & \textbf{Description} \\
  \midrule
  \endhead
  POST & /admin/auth/login                    & PUBLIC     & Admin dashboard login \\
  GET  & /admin/devices                       & Admin JWT  & List all enrolled devices \\
  GET  & /admin/devices/\{deviceId\}          & Admin JWT  & Device detail + event timeline \\
  PUT  & /admin/devices/\{deviceId\}/block    & Admin JWT  & Block/unblock ERP access \\
  GET  & /admin/threats                       & Admin JWT  & Recent threat events (paginated) \\
  PUT  & /admin/threats/\{eventId\}/acknowledge & Admin JWT & Mark alert as acknowledged \\
  GET  & /admin/analytics/summary             & Admin JWT  & Dashboard summary stats \\
  GET  & /admin/analytics/trends              & Admin JWT  & Time-series data for charts \\
  GET  & /admin/patient-zero                  & Admin JWT  & Active threat spread clusters \\
  PUT  & /admin/policy                        & Admin JWT  & Update global access policy \\
  GET  & /admin/reports/export                & Admin JWT  & Export report as PDF/CSV \\
  \bottomrule
  \caption{Admin REST Endpoints}
\end{longtable}

\section{WebSocket (STOMP over SockJS)}

\noindent WebSocket endpoint: \texttt{/ws/admin}

\begin{longtable}{@{}p{5cm}p{9cm}@{}}
  \toprule
  \textbf{STOMP Topic} & \textbf{Description} \\
  \midrule
  \endhead
  /topic/alerts       & New threat alert — broadcasts to all admin clients \\
  /topic/scores       & Device score update event \\
  /topic/devices      & Device online/offline status change \\
  /topic/patient-zero & New cluster detected \\
  \bottomrule
  \caption{WebSocket STOMP Topics}
\end{longtable}

%% ============================================================
%%  CHAPTER 10 — RABBITMQ QUEUE DESIGN
%% ============================================================
\chapter{RabbitMQ Queue Design}

Aegis uses a \textbf{topic exchange} named \texttt{aegis.events}. All queues have associated dead-letter queues (DLQ). This design decouples event ingestion from processing and handles the 9am burst of 200+ simultaneous device connections without database contention.

\begin{longtable}{@{}p{4.5cm}p{2.5cm}p{3cm}p{4cm}@{}}
  \toprule
  \textbf{Queue} & \textbf{Routing Key} & \textbf{Consumer} & \textbf{Purpose} \\
  \midrule
  \endhead
  aegis.telemetry.events & telemetry.\# & TelemetryConsumer    & Heartbeats; updates score in DB and Redis; enforces ERP block \\
  aegis.threat.alerts    & threat.\#    & ThreatAlertConsumer  & Threat reports; rule engine; WebSocket push; Patient Zero \\
  aegis.score.updates    & score.\#     & ScoreUpdateConsumer  & Score recalculation after threat; WebSocket broadcast \\
  \bottomrule
  \caption{RabbitMQ Queue Design}
\end{longtable}

%% ============================================================
%%  CHAPTER 11 — SECURITY AND PRIVACY DESIGN
%% ============================================================
\chapter{Security and Privacy Design}

\section{Security Design}

\begin{description}[leftmargin=*, noitemsep]
  \item[Authentication:] Agents use JWT (HS256, 24h expiry) with refresh token rotation. Admin uses JWT (8h expiry) stored in \texttt{httpOnly} cookies (not localStorage). Separate signing secrets for agent tokens vs.\ admin tokens. JWT revocation via Redis blocklist on logout or device unenrollment.

  \item[Transport Security:] All API calls over HTTPS—nginx handles SSL termination. Certificate pinning on both Android and PC agents via OkHttp \texttt{CertificatePinner}. HSTS header enforced. WebSocket over WSS (secure WebSocket).

  \item[Credential Storage:] Android Keystore with AES-256-GCM (hardware-backed where available). PC agent uses OS keychain (Windows Credential Manager / libsecret on Linux). Server uses bcrypt (strength 12) for admin passwords. Zero plaintext credentials ever logged or transmitted.

  \item[Server Hardening:] nginx exposes only ports 80 and 443 externally. All other services live on an internal Docker bridge network. CORS restricted to dashboard origin only. Input validation via \texttt{@Valid} on all DTOs. JSONB payload size limit: 4KB per event. Rate limiting: 60 requests/minute per device UUID.
\end{description}

\section{Privacy Design (DPDP Act 2023 Compliance)}

India's Digital Personal Data Protection Act 2023 requires informed consent for processing personal data. Aegis is designed for compliance through pseudonymous IDs, behaviour-only monitoring, explicit consent, and data minimization.

\textbf{What Aegis collects:}
\begin{itemize}[leftmargin=*, noitemsep]
  \item Device UUID (pseudonymous)
  \item Security score and breakdown (root status, developer options, screen lock)
  \item App permission flags (which dangerous permissions are granted, not content)
  \item Wi-Fi BSSID (AP identifier, not browsing activity)
  \item Network connection endpoints (IPs and ports, not payload)
  \item File system change events (path and hash, not file content)
  \item Process names and PIDs (not command arguments containing user data)
\end{itemize}

\textbf{What Aegis never collects:}
\begin{itemize}[leftmargin=*, noitemsep]
  \item SMS or call content
  \item Browser history or browsing activity
  \item Photos, camera, or microphone access
  \item GPS location coordinates
  \item File content (only hashes for integrity verification)
  \item Keystrokes or user input
\end{itemize}

The first screen on enrollment displays a clear summary of data collected and obtains explicit acceptance. The \texttt{consent\_accepted} flag in the devices table must be \texttt{true} for registration to succeed.

%% ============================================================
%%  CHAPTER 12 — TEST PLAN
%% ============================================================
\chapter{Test Plan}

\textbf{Testing Approach:} Integration, end-to-end, and manual testing across varying network and security conditions.

\textbf{Tools Used:} JUnit (backend), Robolectric (Android), Jest + React Testing Library (frontend), Arduino Serial Monitor, real-world device setups.

\begin{longtable}{@{}p{0.5cm}p{5cm}p{4.5cm}p{2.5cm}p{1.5cm}@{}}
  \toprule
  \textbf{No.} & \textbf{Scenario} & \textbf{Expected Output} & \textbf{Type} & \textbf{Status} \\
  \midrule
  \endhead
  1  & Evil Twin AP detected (BSSID mismatch)         & CRITICAL alert in dashboard within 5s   & Manual       & \\
  2  & Device rooted, score recalculated               & Score drops by 30; ERP access denied    & Manual       & \\
  3  & Developer options enabled                       & Score drops by 20                       & Manual       & \\
  4  & APK installed from unknown source               & HIGH alert; score deduction applied     & Manual       & \\
  5  & FIM: modify file in watched directory           & MEDIUM alert with old/new hash          & Manual       & \\
  6  & FIM: delete system file                         & MEDIUM alert with DELETE event type     & Manual       & \\
  7  & Suspicious process spawned (netcat)             & MEDIUM alert with PID and cmdLine       & Manual       & \\
  8  & Outbound connection to malicious IP             & HIGH alert; IP and AbuseIPDB score      & Manual       & \\
  9  & Port scan behavior (>10 dest IPs/60s)           & HIGH alert; uniqueDestIPs count         & Manual       & \\
  10 & Sophos auto-login on Wi-Fi connect              & HTTP 200 from 192.168.1.8:8090         & Manual       & \\
  11 & ERP access with score $\geq$ 60                 & WebView loads ERP login page            & Manual       & \\
  12 & ERP access with score $<$ 60                    & Access denied screen with reason        & Manual       & \\
  13 & Admin blocks device manually                    & Agent receives block on next heartbeat  & Manual       & \\
  14 & Network offline during threat report            & Event stored locally; retried with backoff & Integration & \\
  15 & 50 simulated agents sending heartbeats          & All processed; no DB contention         & Load Test    & \\
  16 & WebSocket push on new CRITICAL alert            & Dashboard updates within 5s             & Integration  & \\
  17 & Patient Zero: same threat from 3+ devices       & Cluster identified; first device flagged & Integration & \\
  \bottomrule
  \caption{Aegis System Test Plan}
\end{longtable}

%% ============================================================
%%  CHAPTER 13 — RESULTS
%% ============================================================
\chapter{Results}

\begin{figure}[H]
  \centering
  \fbox{\parbox{0.9\linewidth}{\centering\vspace{80pt}
    \textit{[Insert screenshot: Admin Dashboard — Device Health Board with score badges, risk levels, and filter controls]}
    \vspace{80pt}}}
  \caption{Admin Dashboard — Device Health Board}
  \label{fig:dashboard}
\end{figure}

\begin{figure}[H]
  \centering
  \fbox{\parbox{0.9\linewidth}{\centering\vspace{80pt}
    \textit{[Insert screenshot: Live Alert Feed — CRITICAL Evil Twin alert card with device name, BSSID, and timestamp]}
    \vspace{80pt}}}
  \caption{Live Alert Feed — CRITICAL Threat Alert}
  \label{fig:alertfeed}
\end{figure}

\begin{figure}[H]
  \centering
  \fbox{\parbox{0.9\linewidth}{\centering\vspace{80pt}
    \textit{[Insert screenshot: Android Agent App — Security Score screen showing score breakdown and ERP access status]}
    \vspace{80pt}}}
  \caption{Android Agent — Security Score Screen}
  \label{fig:androidscore}
\end{figure}

\begin{figure}[H]
  \centering
  \fbox{\parbox{0.9\linewidth}{\centering\vspace{80pt}
    \textit{[Insert screenshot: RabbitMQ Management UI at localhost:15672 showing live queue depths during load test]}
    \vspace{80pt}}}
  \caption{RabbitMQ Management UI — Live Queue Depths}
  \label{fig:rabbitmq}
\end{figure}

\noindent Key measured results:
\begin{itemize}[leftmargin=*, noitemsep]
  \item End-to-end latency (event on device $\rightarrow$ dashboard alert): $<$ 5 seconds under normal load.
  \item 50 simulated agents sending concurrent heartbeats: all processed without database contention, validating the RabbitMQ buffering design.
  \item Security Score recalculation and ERP policy enforcement: consistent within one heartbeat cycle ($\leq$ 30 seconds) of a threat event.
  \item Evil Twin detection accuracy: 100\% for known rogue BSSIDs in test environment.
  \item FIM baseline generation for \texttt{/etc} directory (approx.\ 3000 files): completed in $<$ 8 seconds on a lab PC.
\end{itemize}

%% ============================================================
%%  CHAPTER 14 — CONCLUSION
%% ============================================================
\chapter{Conclusion}

\section{Summary}

Aegis successfully implements a multi-layer Distributed Host-Based Intrusion Detection System tailored to the specific security needs of Walchand College of Engineering. By deploying lightweight Android and PC agents alongside a robust Spring Boot backend and real-time React dashboard, the system provides endpoint-level visibility that Sophos's perimeter-only firewall cannot offer.

The system's decoupled architecture—with RabbitMQ absorbing burst traffic and Redis caching live device state—ensures both scalability and low latency. The Security Score model and Zero-Trust ERP gating represent a practical, actionable response to the credential theft risks posed by Evil Twin attacks and phishing on campus networks.

\section{Limitations}

\begin{itemize}[leftmargin=*, noitemsep]
  \item Evil Twin detection relies on a hardcoded/server-maintained BSSID whitelist; new legitimate APs must be manually added by administrators.
  \item FIM on high-activity directories may produce false positives for legitimate automated updates; the debounce window mitigates but does not eliminate this.
  \item Network connection scanning relies on parsing \texttt{netstat} output rather than deep packet inspection, limiting payload-level visibility.
  \item macOS PC agent support is deferred to future work.
  \item The system cannot block network connections at the OS level without elevated (root/admin) privileges, which are intentionally avoided for security reasons.
\end{itemize}

\section{Future Scope}

\begin{itemize}[leftmargin=*, noitemsep]
  \item \textbf{Machine Learning anomaly detection:} Replace the static rule engine with a behavioral baseline ML model to detect novel attack patterns.
  \item \textbf{Integration with college AD/LDAP:} Automatic PRN-to-identity mapping without manual PRN entry on enrollment.
  \item \textbf{macOS agent:} Extend PC agent to macOS via LaunchAgent plist and \texttt{.dmg} installer.
  \item \textbf{SIEM integration:} Export events to industry-standard SIEM platforms (Wazuh, Elastic) for organizations that already have security infrastructure.
  \item \textbf{SaaS model:} Multi-tenant deployment targeting SMEs that lack budget for enterprise-grade SIEM solutions.
  \item \textbf{Mobile dashboard:} React Native or Flutter admin app for on-the-go incident response.
  \item \textbf{OTP interception detection:} Detect apps requesting SMS permissions that could intercept ERP OTP tokens.
\end{itemize}

%% ============================================================
%%  CHAPTER 15 — SCHOLARLY CONTRIBUTION
%% ============================================================
\chapter{Scholarly Contribution}

The project \textit{Aegis: Distributed Host-Based Intrusion Detection System} was developed as part of the TY CSE Mini Project II at Walchand College of Engineering, Sangli.

It contributes to the domains of \textbf{Cyber Security \& Threat Intelligence}, \textbf{Digital Forensics \& Incident Response}, \textbf{Distributed Systems \& Message Brokering}, and \textbf{Low-Level System Design}.

The project demonstrates practical implementation of:
\begin{itemize}[leftmargin=*, noitemsep]
  \item Distributed endpoint monitoring across heterogeneous device types (Android and PC).
  \item Asynchronous event-driven architecture using RabbitMQ for burst-tolerant telemetry ingestion.
  \item Real-time WebSocket push architecture for live security dashboards.
  \item Privacy-preserving security monitoring compliant with DPDP Act 2023.
  \item Zero-Trust enforcement via dynamic Security Score gating on critical infrastructure access.
\end{itemize}

The system design, architecture decisions, API specifications, and testing results are fully documented to serve as a reference for future distributed security systems in campus or SME environments.

The team aims to enhance the system further with machine learning-based anomaly detection and possible submission of findings to technical symposiums or conferences related to campus cybersecurity and distributed intrusion detection.

%% ============================================================
%%  REFERENCES
%% ============================================================
\chapter*{References}
\addcontentsline{toc}{chapter}{References}

\begin{enumerate}[leftmargin=*, noitemsep]
  \item Sysdig (2024). \textit{What is HIDS (Host-Based Intrusion Detection System)?}. \url{https://www.sysdig.com/learn-cloud-native/what-is-hids}

  \item RabbitMQ Documentation (2024). \textit{RabbitMQ: One broker to queue them all}. \url{https://www.rabbitmq.com/}

  \item Wazuh (2023). \textit{Migrating from OSSEC to Wazuh: Open Source Security Platform}. \url{https://wazuh.com/blog/migrating-from-ossec-to-wazuh/}

  \item Tripwire (2022). \textit{File Integrity Monitoring (FIM): What It Is \& How It Works}. \url{https://www.tripwire.com/state-of-security/file-integrity-monitoring}

  \item International Journal of Computer Applications (2013). \textit{Host-based Intrusion Detection and Prevention System (HIDPS)}. \url{https://research.ijcaonline.org/volume69/number26/pxc3888419.pdf}

  \item Spring Framework Documentation (2024). \textit{Spring Boot 3 Reference Guide}. \url{https://docs.spring.io/spring-boot/docs/current/reference/html/}

  \item Android Developers (2024). \textit{WorkManager: Schedule tasks with WorkManager}. \url{https://developer.android.com/topic/libraries/architecture/workmanager}

  \item Android Developers (2024). \textit{Android Keystore System}. \url{https://developer.android.com/privacy-and-security/keystore}

  \item Ministry of Electronics and IT, Government of India (2023). \textit{The Digital Personal Data Protection Act, 2023}. \url{https://www.meity.gov.in/writereaddata/files/Digital_Personal_Data_Protection_Act_2023.pdf}

  \item AbuseIPDB (2024). \textit{AbuseIPDB: IP Address Abuse Reports}. \url{https://www.abuseipdb.com/}

  \item Leaflet (2024). \textit{Leaflet: an open-source JavaScript library for mobile-friendly interactive maps}. \url{https://leafletjs.com/}

  \item Flyway (2024). \textit{Flyway by Redgate: Database Migrations Made Easy}. \url{https://flywaydb.org/}
\end{enumerate}

\end{document}